\cotefile{34}
\scanpage{196}
\mspage{II.44}
% archive sheet number: 195
% lines 1--12
\deleted{\unsure{Et|Si} $X$ est en car $p>0$, alors $X$ est \unsure{gros}.}

Si $X$ est en car $p>0$ $(p=0)$, alors $X\to\operatorname{Spec}\mathbb{F}_p=S$ induit \deleted{\unsure{un}} \unsure{quasi-isom.} $\pi_1(X,\unsure{\xi|\eta})\unsure{-|\to}\pi_1(S,\xi)\simeq\widehat{\mathbb{Z}}$, donc le $\pi_1(X,\xi)^{\unsure{\mathrm{alg}}}\to(\widehat{\mathbb{Z}}^{\mathrm{alg}})$ est un quasi-\deleted{\unsure{isomorphisme}}\inserted{\unsure{épim.}}, et $\widehat{\mathbb{Z}}^{\mathrm{alg}}$ (comme on sait \unsure{aussitôt}) est de dim infinie. Si $X$ est en car $p>0$, alors \unsure{faisant} $S=\operatorname{Spec}\mathbb{Z}[1/\unsure{\ell|l}]$, \deleted{\unsure{on} $\Pi=\pi_1(X,\xi)\to\pi_1(S,$} en raisonnant de $=$, on montre que $\pi_1(S)$ a \deleted{\unsure{un quotient}} comme plus gd \unsure{quotient} abélien $\mathbb{Z}_\ell^*$, qui est quasi-isomorphe à $\mathbb{Z}_\ell$, or $\mathbb{Z}_\ell^{\mathrm{alg}}$ est aussi de dim. infinie ....

% lines 13--21
6.16. Revenons aux conditions générales de 6.13. L'intérêt de la considération de \deleted{\unsure{\(\Pi\)}} $\Pi=\pi_1^{\mathrm{arith}}(X,\xi)$ tient au fait que tous les faisceaux $\ell$-adiques \unsure{constants tordus} qui \unsure{rencontrons|se rencontrent} pratiquement sont de nature arithmétique; \unsure{i.e.|ou.} proviennent déjà d'un tel faisceau sur un $X_\lambda$, \unsure{est unique modulo isom.|et unique modulo isom.} \unsure{?}, sur un $X_\mu$, avec $\mu\gg\lambda$. \deleted{En d'autres termes,} De $=$ pour les \unsure{hom} de faisceaux $\ell$-adiques. En d'autres termes,
